% =====================================================================
%  UAM-TOPOMAT : equations formelles certifiees par SymPy
%  Universite Amadou Mahtar Mbow / UFR STA / Cellule Physique
%  Realise par : Ndiogou Ndiaye -- Annee universitaire 2026-2027
%  Fichier genere automatiquement par src/quantum/symbolic_proofs.py
% =====================================================================
\documentclass[11pt,a4paper,french]{article}
\usepackage[utf8]{inputenc}
\usepackage[T1]{fontenc}
\usepackage[french]{babel}
\usepackage{amsmath,amssymb,amsthm}
\usepackage{booktabs}
\usepackage[a4paper,margin=2.5cm]{geometry}
\usepackage{tcolorbox}
\newcommand{\UAM}{Universite Amadou Mahtar Mbow}
\newcommand{\UFR}{UFR Sciences et Technologies Avancees}

\title{\textbf{UAM-TOPOMAT}\\[0.4em] \large Certificats formels des six equations fondatrices}
\author{Ndiogou Ndiaye}
\date{\UAM \quad---\quad \UFR \quad---\quad Cellule Physique\\Annee universitaire 2026--2027}

\begin{document}
\maketitle

\begin{tcolorbox}[colback=blue!5,colframe=blue!50!black,title=Certification]
Les six theoremes ci-dessous sont derives symboliquement avec SymPy (Python pur, sans SciPy). Tous les residus symboliques sont nuls et les equations LaTeX sont generees par \texttt{sympy.latex}.
\end{tcolorbox}

\tableofcontents

\section{1. Dispersion Tight-Binding et vitesse de groupe d'Hellmann--Feynman}

Pour un Hamiltonien hermitien $H(\mathbf{k})$ dépendant du vecteur d'onde, le théorème d'Hellmann--Feynman donne $\partial_k E_n=\langle u_{n\mathbf{k}}|\partial_k H|u_{n\mathbf{k}}\rangle$, d'où $\mathbf{v}_g^{(n)}=\hbar^{-1}\partial_{\mathbf{k}}E_n$. La chaîne dimérisée $H=d_x\sigma_x+d_y\sigma_y$ est traitée exactement.

\begin{align}
H_{ij}(\mathbf{k}) = \sum_{\mathbf{R}} t_{ij}(\mathbf{R})\,e^{i\mathbf{k}\cdot\mathbf{R}} \\
E_{\pm} = \sqrt{t_{2}^{2} \sin^{2}{\left(a k \right)} + \left(- t_{1} - t_{2} \cos{\left(a k \right)}\right)^{2}} \\
v_g^{(n)} = \frac{a t_{2}^{2} \sin{\left(a k \right)} \cos{\left(a k \right)} + a t_{2} \left(- t_{1} - t_{2} \cos{\left(a k \right)}\right) \sin{\left(a k \right)}}{\hbar \sqrt{t_{2}^{2} \sin^{2}{\left(a k \right)} + \left(- t_{1} - t_{2} \cos{\left(a k \right)}\right)^{2}}} \\
\varepsilon+2t\cos(ak) = \varepsilon + 2 t \cos{\left(a k \right)} \\
v_g^{\mathrm{1D}} = - \frac{2 a t \sin{\left(a k \right)}}{\hbar}
\end{align}

\noindent\emph{Vérification formelle :} les résidus SymPy suivants sont nuls :
\begin{align}
R_{(\mathrm{Hellmann\, Feynman})} = 0 \\
R_{(\mathrm{groupe\, 1D})} = 0
\end{align}

\section{2. Courbure de Berry et invariant topologique de Chern}

Pour un modèle à deux bandes $H=\mathbf{d}\cdot\boldsymbol{\sigma}$, la courbure de la bande inférieure vaut $\Omega_-=-\tfrac12\hat{\mathbf{d}}\cdot(\partial_{k_x}\hat{\mathbf{d}}\times\partial_{k_y}\hat{\mathbf{d}})$, et $C_n=(2\pi)^{-1}\int_{\mathrm{BZ}}\Omega_n\,d^2k\in\mathbb{Z}$. Le cône de Dirac massif contribue $-\tfrac12\operatorname{sgn}(m)$ ; les deux cônes du modèle de Haldane donnent $C=\pm1$ ou $0$.

\begin{align}
\Omega_n(\mathbf{k}) = \partial_{k_x}A_y-\partial_{k_y}A_x \\
\hat{\mathbf{d}}\cdot(\partial_{k_x}\hat{\mathbf{d}}\times\partial_{k_y}\hat{\mathbf{d}}) = \left(\frac{\partial}{\partial k_{x}} \theta{\left(k_{x},k_{y} \right)} \frac{\partial}{\partial k_{y}} \varphi{\left(k_{x},k_{y} \right)} - \frac{\partial}{\partial k_{y}} \theta{\left(k_{x},k_{y} \right)} \frac{\partial}{\partial k_{x}} \varphi{\left(k_{x},k_{y} \right)}\right) \sin{\left(\theta{\left(k_{x},k_{y} \right)} \right)} \\
C_n = \frac{1}{2\pi}\int_{\mathrm{BZ}}\Omega_n(\mathbf{k})\,d^2k\in\mathbb{Z} \\
\Omega_-^{\mathrm{Dirac}} = - \frac{m v^{2}}{2 \left(k^{2} v^{2} + m^{2}\right)^{\frac{3}{2}}} \\
\frac{1}{2\pi}\int\Omega_-^{\mathrm{Dirac}}\,d^2k = - \frac{1}{2} \\
C_{\mathrm{Haldane}} = - \frac{\operatorname{sign}{\left(M - M_{c} \right)}}{2} + \frac{\operatorname{sign}{\left(M + M_{c} \right)}}{2}
\end{align}

\noindent\emph{Vérification formelle :} les résidus SymPy suivants sont nuls :
\begin{align}
R_{(\mathrm{identite\, geometrique})} = 0 \\
R_{(\mathrm{exemple\, theta\, egal\, kx})} = 0 \\
R_{(\mathrm{c\, haldane\, plus\, un})} = 0 \\
R_{(\mathrm{c\, haldane\, moins\, un})} = 0 \\
R_{(\mathrm{c\, haldane\, trivial})} = 0 \\
R_{(\mathrm{cft\, half})} = 0
\end{align}

\section{3. Rendement limite de Shockley--Queisser}

Sous bilan détaillé, $J(V)=J_{sc}-J_0(e^{qV/k_BT}-1)$ avec $J_0=q(2\pi/h^3c^2)\int_{E_g}^{\infty}E^2/(e^{E/k_BT}-1)\,dE$. La condition $\partial_V(VJ)=0$ équivaut exactement à $(1+v)e^v=1+J_{sc}/J_0$, avec $v=qV_{mp}/k_BT$.

\begin{align}
J(V) = - J_0 \left(e^{\frac{V q}{T k_B}} - 1\right) + J_{sc} \\
J_0 = \frac{2 \pi E^{2} \int_{E_g}^{\infty} q}{c^{2} h^{3} \left(e^{\frac{E}{T k_B}} - 1\right)} \\
(1+v)e^v = J_{sc}/J_0 + 1 \\
V_{mp} = \frac{T k_B v}{q} \\
J_{mp} = J_0 v e^{v} \\
\eta_{\mathrm{SQ}} = \frac{J_{mp} V_{mp}}{P_{\mathrm{in}}}
\end{align}

\noindent\emph{Vérification formelle :} les résidus SymPy suivants sont nuls :
\begin{align}
R_{(\mathrm{point\, puissance\, max})} = 0 \\
R_{(\mathrm{courant\, optimal})} = 0
\end{align}

\section{4. Bilan radiatif--convectif de Stefan--Boltzmann}

La constante de Stefan--Boltzmann dérive du rayonnement du corps noir : $\sigma_{SB}=2\pi^5k_B^4/(15h^3c^2)$. À l'équilibre stationnaire, $f(T_j)=\alpha G-\eta(T_j)G-h_{conv}(T_j-T_{amb})-\varepsilon\sigma_{SB}(T_j^4-T_{sky}^4)=0$ avec $T_{sky}=T_{amb}-8$~K. Comme $f'(T_j)=-\beta\eta(T_j)G-h_{conv}-4\varepsilon\sigma_{SB}T_j^3<0$, la racine est unique.

\begin{align}
\sigma_{\mathrm{SB}} = 1.85294433695108 \cdot 10^{-10} \pi^{5} \\
\sigma_{\mathrm{SB}} = 5.670374419 \cdot 10^{-8} \\
P_{\mathrm{abs}} = \alpha G \\
P_{\mathrm{rad}} = \varepsilon \sigma_{\mathrm{SB}} \left(T_j^{4} - \left(T_{amb} - \Delta T\right)^{4}\right) \\
f{\left(T_j \right)} = \alpha G - G \eta(T_j) - \varepsilon \sigma_{\mathrm{SB}} \left(T_j^{4} - \left(T_{amb} - \Delta T\right)^{4}\right) - h_{\mathrm{conv}} \left(- T_{amb} + T_j\right) \\
f'(T_j) = - G \beta \eta(T_j) - 4 T_j^{3} \varepsilon \sigma_{\mathrm{SB}} - h_{\mathrm{conv}}
\end{align}

\noindent\emph{Vérification formelle :} les résidus SymPy suivants sont nuls :
\begin{align}
R_{(\mathrm{derivee\, bilan})} = 0 \\
R_{(\mathrm{constante\, sigma})} = 0 \\
R_{(\mathrm{monotonie})} = 0
\end{align}

\section{5. Scharfetter--Gummel conservatif et relation d'Einstein}

Sur un intervalle de pas $h$ à potentiel linéaire, l'intégration exacte de la dérive--diffusion donne le schéma exponentiel $J_{n,i+1/2}=q\mu_nV_T/h\,[B(\Delta u)n_{i+1}-B(-\Delta u)n_i]$ avec $B(z)=z/(e^z-1)$ et $\Delta u=(u_{i+1}-u_i)$. En variables de Slotboom $w=ne^{-\phi/V_T}$, le flux devient $A_i(w_{i+1}-w_i)$ avec $A_i>0$, ce qui garantit la conservation exacte du courant. Le flux net à l'équilibre s'annule si et seulement si $D=\mu V_T$ (Einstein).

\begin{align}
B(z) = \frac{z}{e^{z} - 1} \\
B(-z) = z + \frac{z}{e^{z} - 1} \\
\frac{z}{e^{z} - 1} = \frac{z^{2}}{12} - \frac{z}{2} + 1 \\
D = V_T \mu \\
V_T = \frac{T_j k_B}{q} \\
J_{n,i+1/2} = \frac{V_T \mu q \left(- B(-\Delta u) n_i + B(\Delta u) n_{i+1}\right)}{h}
\end{align}

\noindent\emph{Vérification formelle :} les résidus SymPy suivants sont nuls :
\begin{align}
R_{(\mathrm{identite\, bernoulli})} = 0 \\
R_{(\mathrm{identite\, exp})} = 0 \\
R_{(\mathrm{slotboom})} = 0 \\
R_{(\mathrm{developpement\, bernoulli})} = 0 \\
R_{(\mathrm{einstein})} = 0
\end{align}

\section{6. Capacité quantique et double couche de Helmholtz en série}

La capacité quantique d'une électrode MXène est $C_Q=e^2\,\mathrm{DOS}(E_F+eV)/A$. En série avec la capacité de Helmholtz $C_H(T)$, $C_{tot}^{-1}=C_Q^{-1}+C_H^{-1}$, donc $C_{tot}=C_QC_H/(C_Q+C_H)<\min(C_Q,C_H)$ et $\partial C_{tot}/\partial C_Q=C_H^2/(C_Q+C_H)^2>0$.

\begin{align}
C_Q = \frac{\mathrm{DOS}(E_F+eV) e^{2}}{A} \\
C_{\mathrm{tot}}^{-1} = \frac{1}{C_Q} + \frac{1}{C_H} \\
C_{\mathrm{tot}} = \frac{C_H C_Q}{C_H + C_Q} \\
\partial C_{\mathrm{tot}}/\partial C_Q = \frac{C_H^{2}}{\left(C_H + C_Q\right)^{2}} \\
q^2/(2C_{\mathrm{tot}}) = \frac{q^{2} \left(C_H + C_Q\right)}{2 C_H C_Q}
\end{align}

\noindent\emph{Vérification formelle :} les résidus SymPy suivants sont nuls :
\begin{align}
R_{(\mathrm{loi\, serie})} = 0 \\
R_{(\mathrm{equilibre\, charge})} = 0 \\
R_{(\mathrm{derivee})} = 0 \\
R_{(\mathrm{energie})} = 0 \\
R_{(\mathrm{borne\, superieure})} = 0
\end{align}

\section{Verification numerique des champions UAM}

Les deux champions certifies satisfont l'ensemble des bornes theoriques exactes etant donnees par les six theoremes.

\begin{table}[htbp]
\centering
\footnotesize
\begin{tabular}{llrr}
\toprule
Champion & Borne & Valeur & Statut \\
\midrule
UAM-1384 & 0.5 <= E\_g <= 2.8 eV & 1.38972 & OK \\
UAM-1384 & mu\_n > mu\_p > 0 & 2.3666 & OK \\
UAM-1384 & C entier, |C| <= 1 & 1 & OK \\
UAM-1384 & eta <= eta\_SQ(E\_g, T\_ref) & 0.228413 & OK \\
UAM-1384 & eta <= eta\_SQ\textasciicircum{}max & 0.228413 & OK \\
UAM-1384 & q V\_oc <= E\_g & 1.07358 & OK \\
UAM-1384 & 0 < FF < 1 & 0.695124 & OK \\
UAM-1384 & J\_sc > 0 & 32.1374 & OK \\
UAM-1384 & T\_j > T\_amb & 349.071 & OK \\
UAM-1384 & |T\_j - T\_j\_SB| < 0.1 K & 0.00766679 & OK \\
UAM-1384 & f(T\_j) = 0 a 1e-6 W/m\textasciicircum{}2 pres & 2.27374e-13 & OK \\
UAM-1384 & h\_conv (T\_j - T\_amb) < G (1 - eta) & 362.666 & OK \\
UAM-1384 & 0.02 < k\_B T\_j/q < 0.05 V & 0.0300806 & OK \\
UAM-1384 & D\_n/D\_p = mu\_n/mu\_p & 2.3666 & OK \\
UAM-1384 & beta\_T < 0 & -0.00191918 & OK \\
UAM-1384 & 0 < R\_47 <= 100 \% (C\_tot < min(C\_Q, C\_H)) & 98.1051 & OK \\
UAM-1384 & 0 < stabilite <= 100 \% & 95.8657 & OK \\
UAM-1384 & 0 < E\_sp < 100 Wh/kg & 3.42148 & OK \\
UAM-1517 & 0.5 <= E\_g <= 2.8 eV & 1.86499 & OK \\
UAM-1517 & mu\_n > mu\_p > 0 & 2.15096 & OK \\
UAM-1517 & C entier, |C| <= 1 & -1 & OK \\
UAM-1517 & eta <= eta\_SQ(E\_g, T\_ref) & 0.180047 & OK \\
UAM-1517 & eta <= eta\_SQ\textasciicircum{}max & 0.180047 & OK \\
UAM-1517 & q V\_oc <= E\_g & 1.51517 & OK \\
UAM-1517 & 0 < FF < 1 & 0.713749 & OK \\
UAM-1517 & J\_sc > 0 & 17.4811 & OK \\
UAM-1517 & T\_j > T\_amb & 347.243 & OK \\
UAM-1517 & |T\_j - T\_j\_SB| < 0.1 K & 0.0119542 & OK \\
UAM-1517 & f(T\_j) = 0 a 1e-6 W/m\textasciicircum{}2 pres & -5.68434e-13 & OK \\
UAM-1517 & h\_conv (T\_j - T\_amb) < G (1 - eta) & 339.745 & OK \\
UAM-1517 & 0.02 < k\_B T\_j/q < 0.05 V & 0.0299231 & OK \\
UAM-1517 & D\_n/D\_p = mu\_n/mu\_p & 2.15096 & OK \\
UAM-1517 & beta\_T < 0 & -0.00174345 & OK \\
UAM-1517 & 0 < R\_47 <= 100 \% (C\_tot < min(C\_Q, C\_H)) & 98.1051 & OK \\
UAM-1517 & 0 < stabilite <= 100 \% & 98.1051 & OK \\
UAM-1517 & 0 < E\_sp < 100 Wh/kg & 5.97398 & OK \\
\bottomrule
\end{tabular}
\caption{Verification des bornes theoriques exactes pour UAM-1384 et UAM-1517.}
\end{table}

\begin{tcolorbox}[colback=green!5,colframe=green!50!black,title=Conclusion]
Les six preuves symboliques sont exactes et les deux champions respectent toutes les bornes theoriques : le certificant est complet.
\end{tcolorbox}

\end{document}
